\begin{frame}{Crítca ao emprego do método}
\begin{columns}[T,totalwidth=\textwidth]
\column{0.58\textwidth}
\begin{itemize}
    \item Dada a heterogeneidade da aplicação do método Case Survey é procurado as diretrizes do métodos em artigos seminais.
    \item O procedimento clássico envolve: selecionar casos já existentes na literatura, criar um esquema de codificação, aplicar múltiplos avaliadores e analisar estatisticamente o conjunto codificado.
\end{itemize}

\column{0.50\textwidth}
\centering
\begin{tikzpicture}[node distance=0.35cm]
\node[draw=steel, fill=soft, rounded corners=2pt, minimum width=4.5cm, minimum height=0.7cm] (c1) {Casos publicados};
\node[draw=steel, fill=soft, rounded corners=2pt, minimum width=4.5cm, minimum height=0.7cm, below=of c1] (c2) {Esquema de codificação};
\node[draw=steel, fill=soft, rounded corners=2pt, minimum width=4.5cm, minimum height=0.7cm, below=of c2] (c3) {Múltiplos avaliadores};
\node[draw=steel, fill=soft, rounded corners=2pt, minimum width=4.5cm, minimum height=0.7cm, below=of c3] (c4) {Análise estatística};
\foreach \a/\b in {c1/c2,c2/c3,c3/c4}{\draw[-{Latex[length=2mm]}, thick] (\a) -- (\b);}
\end{tikzpicture}
\end{columns}
\end{frame}